\chapter{智能体大语言模型的未来方向}\label{chap:concl}


至此，我们已经来到本书的末尾。下面，我们对全书讨论的内容作一个简要的概括性回顾。首先，回顾有关如何构建智能体大语言模型（agentic LLM）的各章；然后，概括智能体大语言模型的不同应用，以及它们在单智能体与多智能体环境中的行为；最后，在第~\ref{sec:limitations}~节讨论当前方法的局限，以及未来可能的发展。

\section{智能体大语言模型的发展现状}
如第一、二、三部所述，智能体大语言模型的基础在于语言建模（language modeling）、推理、工具使用（tool use），以及不同程度的自主性。

\subsection{技术}
智能体大语言模型是在世界中行动、并根据反馈调整自身行为的推理模型。经过预训练与微调（第~\ref{chap:pipeline}~章）后，它们可以在推理阶段访问外部信息，即使用检索增强生成式人工智能（retrieval-augmented generative AI；第~\ref{sec:rag}~章）。它们也可以在推理阶段向自己提供提示，以验证结果，由此实现自我反思（self-reflection）与上下文强化学习（in-context reinforcement learning），见第~\ref{chap:reasoning}~章。智能体大语言模型已经发展为多模态（multimodal）模型，能够在语言、视觉与动作领域中运行（第~\ref{chap:world}~章）。为了有效地完成这些任务，可以为它们配备预训练的世界模型（world model），使其对物理世界具有初步理解。通过与环境交互，这些世界模型可以在测试阶段得到改进。

这类系统使用的主要机器学习架构是 Transformer 架构。模型通过自监督学习进行预训练，即预测下一个词元；随后，通过监督学习与强化学习，针对特定任务进行微调。

上下文中的思维链（chain of thought，CoT）提示显著提高了模型的推理能力。研究者已经开发出用于微调（例如可验证奖励强化学习，RLVR，以及组相对策略优化，GRPO）与自我反思（例如 ReAct 和 reflexion，见第~\ref{chap:reasoning}~章）的推理阶段方法。围绕智能体与环境的交互，人们也发展出了一系列不同的提示方式（第~\ref{chap:agents}~章）。

\subsection{应用}
第四部介绍了多种便于构建智能体大语言模型的工具（第~\ref{chap:tools}~章）。智能体大语言模型已经被设计用于许多不同应用。在电子商务中，它们可以帮助理解用户偏好，并根据推荐采取行动；在医疗中，它们可以阅读和总结诊断文本；在金融中，它们可以分析财务报告，并为交易员提供建议；在科学研究与教育中，它们可以协助文献研究，生成用于运行实验的脚本，分析结果，帮助撰写论文，以及协助评审论文（第~\ref{sec:assist}~章）。

\subsection{行为}
在第五部，我们从技术视角转向行为视角，研究智能体大语言模型在单智能体与多智能体环境中的表现。在讨论大语言模型的基本行为特征之后，第~\ref{chap:behavior}~章讨论了大语言模型驱动的智能体在两类任务上的表现：一类改编自人类行为研究方法，另一类来自专门为人工智能基准测试而开发的框架。第~\ref{chap:role}~章详细考察了基于预定义角色与协议的多智能体交互。第~\ref{chap:emergent}~章则转向开放式环境，研究智能体之间的小规模交互如何在系统层面涌现出行为结构、规范与惯例。

贯穿第五部所讨论的各种方法与研究，我们发现了若干反复出现的模式。其中包括构念效度（construct validity）问题，即观测结果如何推广到人类情境；提示敏感性（prompt sensitivity）与级联效应，尤其是在多智能体环境中；知行鸿沟（knowing-doing gap），即模型以语言作出预测的能力往往不能很好地转化为行动；结构化推理的积极作用，即当推理过程按照借鉴人类推理的方式加以组织时，增加推理预算尤其有用；以及朝不同方向发展的涌现特性。人们既观察到了期望的方向，例如在相同计算预算下超过基线的表现，也观察到了不期望的方向，例如违反规范。

\section{局限与未来工作}\label{sec:limitations}
智能体大语言模型以大语言模型的强大能力为基础，大幅扩展了后者的能力与用途。不过，它们仍有许多局限，下面我们将逐一讨论。
我们先从记忆、规模扩展、可解释性与认知架构（cognitive architecture）等技术性较强的问题谈起，再转向可能具备的感受能力（sentience）或意识（consciousness），以及这些属性所带来的伦理影响等具有广泛意义的话题。对于其中一些问题，前面的章节已经讨论过可能的解决方向。

\subsection{记忆与自我反思}
语言模型是函数：它们计算得到的输出只取决于输入，各次调用之间不保留记忆。为了在世界中自然地行动、进行对话，并通过自我反思从自身行动中学习，智能体必须能够记忆。对我们而言，身份认同取决于记忆能力。智能体也需要某种形式的记忆。

实现初步记忆的一种简单方式，是把以往的交互历史重新放入提示中；当前的聊天机器人已经成功地采用了这种方法。不过，尽管长上下文模型研究取得了一些进展，上下文窗口的大小仍然有限。智能体记忆方法是一个活跃的研究领域。大语言模型应当学会如何最有效地压缩信息，并将其存入长期或短期记忆。

智能体可以通过检索增强生成式人工智能方法（RAG）访问外部信息，例如从企业数据库中获取信息。
研究者正在开发更复杂的智能体记忆架构，借助向量数据库（vector database），以文本或词元嵌入的形式为智能体提供更长期的记忆（第~\ref{sec:memory}~节）。

\subsection{规模扩展}
随着智能体能够自我反思，并从与环境的交互中学习，我们需要讨论底层语言模型的规模扩展问题。

如何进一步提高性能，是大语言模型发展面临的一项主要挑战。最初，规模是主要决定因素，这也是它们被称为{\em 大}语言模型的原因：只有规模足够大、并在足够大的文本语料库上训练的模型，才能在基准测试中取得足够好的表现；也只有大型模型才展现出上下文学习能力。因此，扩大模型与数据的规模似乎就是前进方向。然而，继续扩展正变得越来越困难，因为互联网上大多数可获得的高质量数据，如今都已经被用于训练现有模型。此外，训练更大的模型、延长训练时间，成本都很高；时间、资金与电力消耗限制了这一方向的发展。现有方法正在进入平台期 \citep{sutskever2024}。

为了进一步提高大语言模型的性能，持续获取具有挑战性的训练数据十分重要（第~\ref{chap:synthetic-data}~章）。一种方法是将智能体大语言模型与环境的交互用作训练数据。如果这些交互数据是全新的、多样的、高质量的，就可以用来进行预训练（另见图~\ref{fig:virtuous}）。请注意，GRPO 和 RLVR 等推理阶段方法已经在利用推理阶段的结果进行微调，并取得了良好效果 \citep{guo2025deepseek}；视觉语言动作模型（vision-language-action model，VLA）也在使用自行生成的动作数据进行训练 \citep{black2024pi_0,brohan2023rt}，见第~\ref{chap:world}~章。

\paragraph{过度思考}
有关思维链的扩展研究也揭示了增加思考投入的局限。思维链在某些领域比在其他领域更有效 \citep{sprague2024cot}。此外，更深入的推理可能导致{\em 过度思考（overthinking）}：投入越多，推理过程反而越低效 \citep{liu2024mind,wu2025more}。相关研究仍在继续，人们正逐步深入理解思维链如何运作 \citep{chen2025does}，以及哪些机制是有效的 \citep{chen2026think}。

\paragraph{开放世界模型与课程学习}
要使规模扩展这一方向奏效，合成训练数据必须具有多样性，并且不同于现有数据。开放世界模型（open world model）的研究，专门以生成多样的训练环境为目标。这类模型的构想是创造无穷无尽的数据集或环境，使智能体能够不断学习、持续提高智能水平 \citep{bruce2024genie}。开放世界模型受到演化算法、质量—多样性（quality-diversity）方法与计算机游戏的启发 \citep{vinyals2019grandmaster,mouret2015illuminating,meta2022human}。

进一步扩展模型性能的另一种方法，是开发更好的训练方法。课程学习（curriculum learning）基于这样一种观察：对大多数学习过程而言，先在简单任务上训练系统，再尝试更困难的任务，效果往往更好。（这既适用于人脑，也适用于人工神经网络。）在大多数领域，都存在由相互关联、层层递进的概念构成的层级结构。（在学校，我们先学加法，再学长除法，甚至更晚才学对数。）研究者正在预训练与后训练中尝试课程学习方法 \citep{chen2025self,elgaar2026curriculum}。

\subsection{幻觉与可解释性}
当前智能体大语言模型的另一项局限是幻觉（hallucination）问题：模型经常给出听起来合理、但事实错误的答案；而且，它们所陈述的推理也未必是所生成输出的因果成因——这种性质被称为\emph{忠实性（faithfulness）}。

可解释人工智能（explainable artificial intelligence，XAI）领域提供了一系列应对这一问题的技术。这些技术从将模型视为黑箱的行为方法与基于归因的方法，例如 SHAP 和 LIME，延伸到基于概念的可解释性（concept-based interpretability），即探查内部表征中人类能够理解的特征。更进一步，第~\ref{chap:mi}~章讨论的机制可解释性（mechanistic interpretability，MI），旨在对 Transformer 模型内部的计算回路进行逆向工程。

稀疏自编码器（sparse autoencoder）、神经透镜（neural lens）与激活修补（activation patching）等工具，已经被用于识别造成幻觉的特定回路，并提供了实证证据，表明在足够大的架构中，思维链推理可以成为模型行为中忠实且具有因果作用的组成部分。随着智能体大语言模型越来越多地部署于高风险场景，机制可解释性为幻觉检测器提供了一条发展路径：当模型内部的置信特征与其外部输出出现偏离时，检测器便会触发。这是构建可问责的智能体系统的关键一步。

再进一步，就是比较人类与 Transformer 的认知架构。语言模型生成自然语言的方式，似乎与人类相似。这是因为生成话语的架构相似，还是出于巧合，抑或另有解释？第~\ref{chap:cognition}~章探讨了这一问题。

\subsection{情绪行为与意识}
人类具有拟人化（anthropomorphism）的倾向：当周围的无生命物体表现得与我们相似时，我们会认为它们的行为具有目的性；有时，即使它们的行为与我们不同，我们也会这样做。我们还可能在不知不觉中将意识赋予这类物体。于是，我们会相信并说出这样的话：{\em 计算机在思考}、{\em 计算机想要某样东西}、{\em 计算机一定是不舒服了}，甚至说{\em 这本书喜欢被阅读}。大语言模型也遭遇了同样的命运，它们的话语有时甚至会让最理性的哲学家相信它们具有目的（然而，它们显然不可能具有目的，因为它们不过是在机械装置上执行的算法；人类则不同，因为我们是有机体）。\footnote{需要指出，鸭子类型（duck typing）的观点并不同意这一点。这种行为主义观点认为，如果它走起来像鸭子，叫起来也像鸭子，我们就把它当作鸭子。}

已经有报告讨论了聊天机器人与人类之间的情感关系；这些关系有时会导致危险行为，甚至更严重的后果。

人工智能哲学家正在积极争论意识问题。人们通常认为，意识意味着能够觉察自我与环境，并具有主观体验、思想和感受 \citep{velmans2009understanding,butlin2023consciousness}。但是，我们该如何检验意识？
一种务实的观点是，对我们的社会而言，真正重要的是行为。如果一个实体采取行动，重要的是行动的后果，而它是否具有意识则没有那么重要。\citet{bengio2025illusions} 认为，即使只是通过行动制造出意识的幻象，也会产生后果；而且，我们将越来越难以区分有意识的实体与无意识的实体。他们主张，无论智能体是在医疗、金融还是科学研究领域行动，社会都必须找到应对之道。功劳与责任最终都必须得到公平归属，对于如何处理自我保护型人工智能系统（self-preserving AI systems），也必须找到答案。

\subsection{社会行为}
我们讨论过大语言模型如何模拟人类推理的例子（第~\ref{chap:behavior}~章），以及它们如何改进多智能体系统中的社会行为模拟。当前智能体社会行为模拟的一项局限是，行为受到提示脚本的强烈影响。虽然有些研究发现了符合现实的社会行为，但许多研究也揭示了大语言模型无法胜任的情形。

此外，大规模研究涌现行为也面临挑战（第~\ref{chap:emergent}~章）。
与所有多智能体模拟一样，验证涌现行为、判断模拟行为有多符合现实，是一项挑战，尽管有些研究报告了令人印象深刻的结果 \citep{park2024generative}。

\subsection{法律与伦理问题}
助手在现实世界中的行动会产生现实后果。当不完美模型的行动会影响现实世界时，能够行动的智能体大语言模型就面临许多法律与伦理挑战。如果金融交易助手产生幻觉，或者自主行驶机器人作出错误推理，那么，在将这些系统部署于现实世界之前，就应当解决责任与法律责任问题。


一个主要研究方向是从技术层面提高安全性。
要确保智能体大语言模型的安全性，就需要超越基于提示的防御，采用集成的多层保障措施。在安全强化学习（safe reinforcement learning）领域，一项重要方向是对智能体的动作施加明确的安全约束，并将这些约束纳入其推理与规划流程 \citep{gu2024review}。

提高安全性的另一种方法，是采用持续对抗训练（continual adversarial training）与自动化红队测试（automated red teaming），增强系统抵御操纵的鲁棒性。其他优先方向还包括开发自我调节与风险感知机制，使智能体能够检测并避免不安全动作，以及建立严格、标准化的安全基准测试 \citep{andriushchenko2024agentharm}。这些措施共同勾勒出一条路线，旨在高风险领域中实现可问责、可信赖的智能体大语言模型部署。

不过，对于谁应当为智能体的行动负责，法律与伦理层面的问题仍然存在。如今，智能体开始在现实世界中投入使用，解决这些问题变得至关重要。
显然，只有解决许多安全、伦理与信任问题，我们才能充分受益于智能体大语言模型的广泛潜能。
安全与伦理将成为智能体大语言模型研究议程中的重要主题，值得以专门的综述与著作加以讨论 \citep{gan2024navigating,jiao2025navigating,raza2025responsible}。




\subsection{犯罪}
当我们考虑犯罪行为时，有关伦理与意识的讨论就变得十分具体。守法的人和违法的人都可以利用技术进步。技术也可能被用于从事不道德行为，而当前尚无有效法律来约束这些行为。
随着智能体大语言模型能够展现出越来越可信的行为，我们将难以分辨真实行为与人工生成的行为。这可能导致危险的伦理情境，例如犯罪分子利用这项技术，影响毫无戒心的受害者的行为。

合法说服、滥用信任与欺诈之间的界线，可能并不清楚。
借助社交网络、营销、间谍活动与网络犯罪中的说服算法，智能体大语言模型将大幅提高社会工程（social engineering）的复杂程度 \citep{kshetri2025transforming,lupinacci2025dark}。尽管智能体大语言模型也可以用于打击犯罪 \citep{challita2025redteamllm,andriushchenko2024agentharm}，但社会对这些风险的认识，还需要一段时间才能跟上技术发展。
尤其是那些不太熟悉社交网络与大语言模型的人，可能容易遭受信任诈骗，以及利用智能体大语言模型实施的网络犯罪。


\subsection{应用的局限}
智能体大语言模型的应用十分广泛。人们已经构建了许多概念验证系统，初看之下，其结果令人印象深刻。然而，深入考察之后，就会发现一些可能不易解决的局限；当智能体被用于现实世界时，这些局限还可能产生严重后果。

\subsubsection{人工智能科研}
许多智能体大语言模型旨在支持科学研究。研究流程的各个阶段都有相应支持：文献总结工具帮助形成研究构想；一些工具可以生成脚本，用于在不同领域运行实验；还有工具支持数据解释、图表制作、写作，以及科学论文评审。的确，产出表面上文笔流畅、呈现精美的研究论文，正变得越来越容易。

在数学领域，有研究报告了令人印象深刻的结果：无论在几何 \citep{trinh2024solving} 还是代数 \citep{hubert2025olympiad} 方面，大语言模型智能体都达到了数学奥林匹克竞赛的水平。这些系统包含 Transformer 网络，但在其他方面是作为数学求解器来设计的，使用 Lean \citep{de2015lean} 等专门的自动定理证明（automated theorem proving）技术，以及强化学习 \citep{silver2017mastering,browne2012survey}。

2025 年，一项基准研究在机器学习的一系列领域中评估了人工智能科研工具 \citep{nathani2025mlgym}。使用当时的前沿模型，作者并未发现这些模型能够生成新颖的假设、算法或架构，也未发现它们能够带来实质性的科学进步。作者特别指出，{\em 现代大语言模型智能体能够成功应对多种多样的定量实验，体现出先进的技能与领域适应能力}；但他们也指出，{\em 目前仍不清楚，科学新颖性这一概念能否成功实现自动化，甚至能否以适合智能体使用的形式进行形式化定义}。



\subsubsection{电子商务}
行动的法律责任，是在电子商务中部署智能体大语言模型的一项主要问题。安全与隐私顾虑会使用户不愿向可能泄露信息的计算机系统提供个人资料与偏好。由此便形成了推荐系统的两难困境：用户不愿详细提供真实偏好，推荐质量就会较低；用户因而不太信任系统，也就更不愿提供信息。只有能够以安全可靠的方式轻松收集高质量偏好信息，并且模型能够给出有力而实用的建议时，公众才会采用电子商务智能体大语言模型。

\subsubsection{交易}
对于金融交易，法律责任至关重要。
有研究报告，大语言模型在总结大量财务报告方面，可以超过人类分析师。然而，金融交易涉及重大利益，错误的交易建议会直接导致经济损失。此外，金融市场具有不可预测性，因此，可能需要一段时间才能识别出一个有问题的模型：对于单笔亏损，往往难以判断它究竟源于模型的缺陷，还是仅仅因为交易运气不佳。模型必须经过充分测试，才会获得足够的信任并投入使用。除非做到万无一失，否则自动交易系统的作用将止于辅助人类决策。

大语言模型在金融领域的成功，大多体现在文本报告分析方面。对于金融数据的时间序列分析，虽然已有一些研究结果发表，但很难获得可靠、可验证的基准测试结果。
尽管有些研究声称取得了良好结果，另一些研究却报告，智能体大语言模型对金融市场结果的预测仍然存疑 \citep{li2026can}。




\subsubsection{医疗}
在医疗领域，要建立可信赖性，结果的准确性，尤其是可解释性，更为重要。模型预测与行动的法律责任和伦理直接相关。对于医疗应用，结果的可解释性至关重要，智能体大语言模型系统将让人类参与决策回路。

与交易领域一样，智能体大语言模型在医疗领域最有前景的成果，也体现在文本分析与总结方面，例如分析和总结会议记录或文献综述。在治疗方案决策、实施手术等其他应用中，成果还比较有限。

医疗很可能是智能体大语言模型最有价值、同时也最为保守的应用领域之一。对于生死攸关的问题，我们必须格外谨慎，而与此同时，潜在收益也十分巨大。



\section{结语}

智能体大语言模型技术正在迅速发展。本书围绕三个相互关联的支柱，梳理了这一发展历程：\emph{推理}、\emph{行动}与\emph{交互}。推理大语言模型通过思维链、自我反思与检索增强，提供认知基础，使模型能够完成多步问题求解、文献综合与复杂决策支持。能够行动的智能体将这一基础延伸到现实世界，利用工具、API 与世界模型，在机器人、金融、医疗等领域协助用户。彼此交互的智能体则更进一步，在多智能体模拟中相互交流，展现出合作、信任与规范形成等涌现社会行为，而这些行为并未在提示中显式编排。

这三个支柱彼此相连。如图~\ref{fig:virtuous}~所示，它们构成一个循环：推理有助于更好地行动与交互，而智能体通过行动和交互生成的数据，又可以用于预训练与微调更好的推理模型。这一良性循环为突破大语言模型规模扩展的平台期提供了一条有前景的路径，因为智能体助手、机器人系统与多智能体模拟在推理阶段产生的交互数据，能够持续补充高质量、多样化的训练信号。

然而，三个支柱都仍然面临重大挑战。幻觉削弱了人们对推理链的信任，并威胁着自我改进系统的稳定性。在高风险领域，行动智能体的安全、法律责任与伦理问责问题仍未解决。大规模多智能体系统中的涌现行为，仍然难以验证和控制。要克服这些挑战，还需要进一步研究机制可解释性、开放世界模型、可扩展的模拟基础设施，以及稳健的安全基准测试。

我们正处于一个早期而又影响深远的时刻。概念验证令人印象深刻，现实世界中的部署正在起步，而关于记忆、元认知、意识与集体智能的理论基础，仍是尚待解答的问题，它们将人工智能与哲学、认知科学和法律联系起来。本书希望为你提供参与研究这些问题所需的技术基础与概念词汇。这一领域需要这样的研究者：既理解当今智能体大语言模型的强大能力，也理解其局限，并能帮助塑造它们的未来。





\bibliographystyle{plainnat}
\bibliography{agenticbook,refs,case,synthetic}

